\chapter{De Giorgi Iteration and Local Boundedness}\label{sec:degiorgi}
%% ========================================================================

% !TEX root = DeGiorgi_Lean_to_Tex.tex

This section is the formal core of the De Giorgi $L^2 \to L^\infty$
argument for non-negative subsolutions of a divergence-form linear
elliptic equation
$\dv(A \nabla u) \le 0$
on the unit ball $\Bz{1} \subset \R^d$, with $A$ measurable, symmetric,
and uniformly elliptic with constants $0 < \lambda \le \Lambda$ (here
$\lambda$ denotes always the global ellipticity constant, never a level
threshold).  The Lean development is organised into six layers, each
realised by a dedicated module:\\
\leanname{DeGiorgiIteration.CutoffAdmissibility},
\leanname{DeGiorgiIteration.Energy},\\
\leanname{DeGiorgiIteration.PreIteration},
\leanname{DeGiorgiIteration.Recurrence}, and
\leanname{DeGiorgiIteration.Linfty},
all reassembled into the public interface
\leanname{DeGiorgiIteration} and reexported by
\leanname{DeGiorgiTheory}.  We follow the same layering below.

Throughout the section $E = \R^d$ with $d \ge 1$, $\Omega \subset E$ is
an open set with finite Lebesgue measure, and $A$ is an elliptic
coefficient field with constants $\lambda$, $\Lambda$ encoded in
\leanname{EllipticCoeff}; we write $\mathrm{ellipticityRatio}(A) =
\Lambda/\lambda$.

\section{Positive-part truncation}

\begin{definition}
\label{def:positivePartSub}
\lean{DeGiorgi.positivePartSub}
\leanok
For any function $u : \alpha \to \R$ and any level $\ell \in \R$ we
define the positive-part truncation
\begin{equation*}
  (u-\ell)_+(x) := \max(u(x) - \ell, 0).
\end{equation*}
\end{definition}

\begin{lemma}
\label{lem:positivePartSub_nonneg}
\lean{DeGiorgi.positivePartSub_nonneg}
\leanok
\uses{def:positivePartSub}
$0 \le (u-\ell)_+(x)$ for every $x$.
\end{lemma}

\begin{proof}
\leanok

Immediate from $\max(\,\cdot\,, 0) \ge 0$.
\end{proof}

\begin{lemma}
\label{lem:positivePartSub_sq_ge_gap_sq_of_lt}
\lean{DeGiorgi.positivePartSub_sq_ge_gap_sq_of_lt}
\leanok
\uses{def:positivePartSub}
If $\theta < \ell$ and $\ell < u(x)$, then
$(\ell - \theta)^2 \le |(u-\theta)_+(x)|^2$.
\end{lemma}

\begin{proof}
\leanok

Under the hypotheses $u(x) - \theta > 0$, hence
$(u-\theta)_+(x) = u(x) - \theta \ge \ell - \theta > 0$.  Squaring this
nonnegative inequality yields the claim.
\end{proof}

\begin{lemma}
\label{lem:positivePartSub_antitone}
\lean{DeGiorgi.positivePartSub_antitone}
\leanok
\uses{def:positivePartSub}
If $\ell_1 \le \ell_2$, then
$(u-\ell_2)_+ \le (u-\ell_1)_+$ pointwise.
\end{lemma}

\begin{proof}
\leanok

$\max(u-\ell_2, 0) \le \max(u-\ell_1, 0)$ since $u-\ell_2 \le u-\ell_1$.
\end{proof}

\begin{lemma}
\label{lem:positivePartSub_le_posPartZero}
\lean{DeGiorgi.positivePartSub_le_posPartZero}
\leanok
\uses{def:positivePartSub}
If $0 \le \ell$, then $(u-\ell)_+ \le u_+ = \max(u,0)$.
\end{lemma}

\begin{proof}
\leanok
\uses{lem:positivePartSub_antitone}
See the Lean formalization.
\end{proof}

\begin{lemma}
\label{lem:le_of_positivePartSub_eq_zero}
\lean{DeGiorgi.le_of_positivePartSub_eq_zero}
\leanok
\uses{def:positivePartSub}
If $(u-\ell)_+(x) = 0$, then $u(x) \le \ell$.
\end{lemma}

\begin{proof}
\leanok

If $u(x) - \ell \ge 0$ then $(u-\ell)_+(x) = u(x) - \ell$, so the
hypothesis forces $u(x) = \ell$.  Otherwise $u(x) < \ell$.
\end{proof}

\section{Cutoff admissibility}

The De Giorgi test functions are obtained by multiplying the
positive-part truncation by the square of a smooth scalar cutoff.

\begin{definition}
\label{def:deGiorgiCutoffTestGeneral}
\lean{DeGiorgi.deGiorgiCutoffTestGeneral}
\leanok
\uses{def:positivePartSub}
Given $\eta, u : E \to \R$ and a level $\ell$,
\begin{equation*}
  \varphi_{\eta, u, \ell}(x) := \eta(x) \cdot \bigl( \eta(x) \cdot
  (u-\ell)_+(x) \bigr) = \eta(x)^2 \, (u-\ell)_+(x).
\end{equation*}
\end{definition}

\begin{lemma}
\label{lem:deGiorgiCutoffTestGeneral_nonneg}
\lean{DeGiorgi.deGiorgiCutoffTestGeneral_nonneg}
\leanok
\uses{def:deGiorgiCutoffTestGeneral}
If $\eta \ge 0$ pointwise, then $\varphi_{\eta, u, \ell} \ge 0$.
\end{lemma}

\begin{proof}
\leanok
\uses{def:positivePartSub, lem:positivePartSub_nonneg}
See the Lean formalization.
\end{proof}

\begin{lemma}
\label{lem:deGiorgiCutoffTestGeneral_hasCompactSupport}
\lean{DeGiorgi.deGiorgiCutoffTestGeneral_hasCompactSupport}
\leanok
\uses{def:deGiorgiCutoffTestGeneral}
If $\eta$ has compact support, so does $\varphi_{\eta,u,\ell}$.
\end{lemma}

\begin{proof}
\leanok
\uses{def:positivePartSub}
See the Lean formalization.
\end{proof}

\begin{lemma}
\label{lem:deGiorgiCutoffTestGeneral_tsupport_subset}
\lean{DeGiorgi.deGiorgiCutoffTestGeneral_tsupport_subset}
\leanok
\uses{def:deGiorgiCutoffTestGeneral}
If $\tsupp\,\eta \subseteq \Omega$, then $\tsupp\,\varphi_{\eta,u,\ell}
\subseteq \Omega$.
\end{lemma}

\begin{proof}
\leanok
\uses{def:positivePartSub}
See the Lean formalization.
\end{proof}

The technical core of this layer is a Sobolev membership statement: if
$u \in W^{1,2}(\Omega)$ and the positive-part truncation
$(u - \ell)_+$ is also in $W^{1,2}$, then for any smooth bounded
cutoff $\eta$ with $\tsupp\,\eta \subseteq \Omega$, the test function
$\varphi_{\eta,u,\ell}$ lies in $W^{1,2}_0(\Omega)$.  In Lean this is
phrased through the witness type \leanname{MemW1pWitness} and the
zero-trace structure \leanname{MemW01p}.

\begin{theorem}
\label{thm:deGiorgiCutoffTest_memW01p_of_truncWitness}
\lean{DeGiorgi.deGiorgiCutoffTest_memW01p_of_truncWitness}
\leanok
\uses{def:MemW1pWitness, def:MemW01p, def:positivePartSub, def:deGiorgiCutoffTestGeneral}
Let $\Omega \subset E$ be open with finite measure.  Let $u$ be such
that $(u-\ell)_+$ admits a $W^{1,2}(\Omega)$ witness.  Let $\eta \in
C^\infty(E)$ satisfy $|\eta| \le C_0$, $\|\nabla \eta\| \le C_1$,
$\eta$ has compact support, and $\tsupp\,\eta \subseteq \Omega$.
Then $\varphi_{\eta,u,\ell} \in \Wop{2}(\Omega)$.
\end{theorem}

\begin{proof}
\leanok
\uses{def:MemW1p, thm:memW01p_of_compactly_supported, lem:MemW1pWitness_mul_smooth_bounded, lem:deGiorgiCutoffTestGeneral_hasCompactSupport, lem:deGiorgiCutoffTestGeneral_tsupport_subset}

By two successive applications of the closure of $W^{1,2}$-witnesses
under multiplication by smooth bounded functions
(\leanname{MemW1pWitness.mul\_smooth\_bounded}), starting from $(u-\ell)_+$
and multiplying twice by $\eta$, one obtains a $W^{1,2}$ witness for
$\eta^2 (u-\ell)_+$.  Since $\varphi_{\eta,u,\ell}$ is compactly
supported with $\tsupp \subseteq \Omega$, the localisation lemma
\leanname{memW01p\_of\_memW1p\_of\_tsupport\_subset} promotes the
$W^{1,2}$-witness to a $W^{1,2}_0(\Omega)$ membership.
\end{proof}

\begin{theorem}
\label{thm:deGiorgiCutoffTest_memW01p_of_posPart}
\lean{DeGiorgi.deGiorgiCutoffTest_memW01p_of_posPart}
\leanok
\uses{def:MemW1pWitness, def:MemW01p, def:deGiorgiCutoffTestGeneral}
Same conclusion, but assuming a $W^{1,2}$-witness for $u$ together
with a generic positive-part closure rule
$\bigl(\, v \in W^{1,2}(\Omega) \implies (v)_+ \in W^{1,2}(\Omega)\,\bigr)$.
\end{theorem}

\begin{proof}
\leanok
\uses{lem:MemW1pWitness_sub_const, def:positivePartSub, thm:deGiorgiCutoffTest_memW01p_of_truncWitness}

Apply the closure rule to $v = u-\ell$ to obtain a witness for
$(u-\ell)_+$, then invoke the previous theorem.
\end{proof}

A concrete version, \leanname{deGiorgiCutoffTest\_memW01p\_of\_posPartApprox},
is obtained by furnishing the $W^{1,2}$-witness for $(u-\ell)_+$ via the
Chapter~02 smooth approximation package for $u-\ell$.  The
ball-restricted variants
\leanname{deGiorgiCutoffTest\_memW01p\_on\_ball\_of\_ballPosPart} and
\leanname{deGiorgiCutoffTest\_memW01p\_on\_ball\_of\_ballPosPartApprox}
specialise these statements to $\Omega = \Ball{s}{x_0}$ and supply
finiteness of the restricted measure automatically.

\begin{definition}
\label{def:positivePartSub_memW1pWitness_on_ball}
\lean{DeGiorgi.positivePartSub_memW1pWitness_on_ball}
\leanok
\uses{def:MemW1pWitness, lem:MemW1pWitness_sub_const, thm:MemW1pWitness_posPart, def:positivePartSub}
On a ball $\Ball{s}{x_0}$, given a $W^{1,2}$-witness for $u$ and a
smooth approximation package for $u - \ell$, the lemma constructs a
$W^{1,2}$-witness for $(u - \ell)_+$ on the ball.
\end{definition}

\section{Weighted Caccioppoli inequality}

We now turn to the analytical heart of the chapter: the weighted
Caccioppoli inequality for truncations of subsolutions.  We use the
abbreviation
\begin{equation*}
  \nabla \eta(x) := (\partial_1 \eta(x), \dots, \partial_d \eta(x))
\end{equation*}
identified with the inner-product representation of the Fréchet
derivative; see \leanname{deGiorgiFderivVec}.  The norm
$\|\nabla \eta(x)\|$ then equals the operator norm
$\|D\eta(x)\|$, which is recorded as
\leanname{deGiorgi\_norm\_fderivVec\_eq\_norm\_fderiv}.

\begin{lemma}
\label{lem:weighted_caccioppoli_pointwise}
\leanok
\noindent\textit{(Formalized as private declaration \leanname{DeGiorgi.weighted\_caccioppoli\_pointwise}.)}\par
Let $\Lambda > 0$, $\eta, \zeta, E, M, v \in \R$, and assume the
pointwise coercivity condition $|M|^2 \le \Lambda E$. Then
\begin{align*}
  2 \eta |v| \zeta |M|
   \le \tfrac{1}{2}\, \eta^2 E + 2\Lambda\, \zeta^2 |v|^2.
\end{align*}
\end{lemma}

\begin{proof}
\leanok

Apply Young's inequality $2ab \le a^2 + b^2$ with $a =
\eta |M| / \sqrt{2\Lambda}$ and $b = \sqrt{2\Lambda}\, \zeta |v|$ to
get
\begin{equation*}
  2 \eta |v| \zeta |M|
   \le \frac{\eta^2 |M|^2}{2\Lambda} + 2\Lambda\, \zeta^2 |v|^2.
\end{equation*}
The hypothesis $|M|^2 \le \Lambda E$ then yields $\eta^2 |M|^2 /
(2\Lambda) \le \tfrac{1}{2} \eta^2 E$, completing the proof.
\end{proof}

\begin{theorem}[Abstract absorption]
\label{thm:weighted_caccioppoli_absorb}
\lean{DeGiorgi.weighted_caccioppoli_absorb}
\leanok
Let $(\alpha, \mu)$ be a measure space and let $E, M, \eta, \zeta, v
: \alpha \to \R$ be measurable functions with $\Lambda > 0$. Assume
the integrability conditions and:
\begin{enumerate}
\item $\int \eta^2 E \, d\mu \le \int 2 \eta |v| \zeta |M| \, d\mu$,
\item $|M|^2 \le \Lambda E$ almost everywhere.
\end{enumerate}
Then $\int \eta^2 E \, d\mu \le 4 \Lambda \int \zeta^2 |v|^2 \, d\mu$.
\end{theorem}

\begin{proof}
\leanok
\uses{lem:weighted_caccioppoli_pointwise}

Apply the pointwise Young inequality almost everywhere, integrate to
get
\begin{equation*}
  \int 2\eta|v|\zeta|M|\, d\mu
   \le \tfrac{1}{2}\int \eta^2 E\, d\mu
        + 2\Lambda \int \zeta^2 |v|^2 \, d\mu,
\end{equation*}
combine with the assumed core estimate to obtain
$\int \eta^2 E\, d\mu \le \tfrac12 \int \eta^2 E\, d\mu + 2\Lambda
\int \zeta^2 |v|^2 \, d\mu$, and absorb to conclude.
\end{proof}

The PDE-facing weighted Caccioppoli inequality is the following
statement, which is the analytical core of the De Giorgi argument.

\begin{theorem}[Weighted Caccioppoli for subsolutions]
\label{thm:caccioppoli_weighted_of_subsolution}
\lean{DeGiorgi.caccioppoli_weighted_of_subsolution}
\leanok
\uses{def:MemW1pWitness, def:MemW01p, def:EllipticCoeff, def:ellipticityRatio, def:IsSubsolution, def:positivePartSub, def:deGiorgiCutoffTestGeneral}
Let $\Omega \subset E$ be open with finite measure, $A$ an elliptic
coefficient field on $\Omega$, $u$ a subsolution
($\leanname{IsSubsolution}\, A\, u$), and $\eta$ a smooth bounded
nonneg function with $|\eta| \le C_0$, $\|\nabla \eta\| \le C_1$,
$\tsupp\,\eta \subseteq \Omega$.  Let $\ell \in \R$ and assume
$(u - \ell)_+ \in W^{1,2}(\Omega)$ and that $\varphi_{\eta,u,\ell}$
admits a $W^{1,2}_0(\Omega)$ structure.  Then
\begin{multline*}
  \int_\Omega \eta(x)^2 \, \bigl\| \nabla(u-\ell)_+(x) \bigr\|^2 \, dx
   \le \\
   4\, \frac{\Lambda}{\lambda}\,
     \int_\Omega \|\nabla \eta(x)\|^2 \,
       \bigl|(u-\ell)_+(x)\bigr|^2 \, dx.
\end{multline*}
\end{theorem}

\begin{proof}
\leanok
\uses{lem:MemW1pWitness_sub_const, lem:MemW1pWitness_mul_smooth_bounded, lem:EllipticCoeff_ae_coercive_nonneg, lem:EllipticCoeff_mulVec_sq_le, lem:EllipticCoeff_quadratic_upper, def:bilinFormIntegrandOfCoeff, lem:aestronglyMeasurable_bilinFormIntegrandOfCoeff, lem:positivePartSub_nonneg, lem:deGiorgiCutoffTestGeneral_nonneg}

Test the subsolution inequality against
$\varphi := \varphi_{\eta,u,\ell} = \eta^2 (u-\ell)_+$, which is
admissible by the cutoff admissibility theorem and is nonnegative.
The subsolution condition gives
$\int_\Omega \langle A\nabla u, \nabla \varphi\rangle \, dx \le 0$.
Computing the gradient of $\varphi$,
\begin{align*}
  \nabla \varphi(x)
   = \eta(x)^2 \, \nabla (u-\ell)_+(x)
     + 2\, \eta(x)\, (u-\ell)_+(x) \, \nabla \eta(x),
\end{align*}
which is the content of
\leanname{deGiorgiCutoffTestWitnessWeighted\_grad}. Combining with the
identification $\nabla u = \nabla(u - \ell)_+$ on $\{u > \ell\}$ and
$\nabla(u-\ell)_+ = 0$ on $\{u \le \ell\}$ (the lemmas
\leanname{truncGrad\_eq\_on\_superlevel} and
\leanname{truncGrad\_zero\_on\_sublevel}), we may rewrite
\begin{multline*}
  0 \ge
   \int_\Omega \eta(x)^2 \,
     \langle A(x) \nabla(u-\ell)_+(x), \nabla(u-\ell)_+(x)\rangle \, dx \\
   + \int_\Omega 2 \eta(x) (u-\ell)_+(x)\,
     \langle A(x) \nabla(u-\ell)_+(x), \nabla \eta(x)\rangle \, dx.
\end{multline*}
By the coercivity bound for $A$ (lemma
$\Lambda$-coercive: $\langle Aw,w\rangle \ge \lambda \|w\|^2$) and the
upper bound (lemma $\langle Aw,w\rangle \le \Lambda\|w\|^2$ together
with $\|Aw\|^2 \le \Lambda \langle Aw,w\rangle$), the first term is
bounded below by $\lambda \int \eta^2 \|\nabla(u-\ell)_+\|^2$, while
the second is dominated in absolute value by
\begin{equation*}
  \int_\Omega 2\, \eta(x) (u-\ell)_+(x) \,
    \|A(x)\nabla(u-\ell)_+(x)\|\,\|\nabla \eta(x)\| \, dx.
\end{equation*}
Setting $E := \langle A\nabla(u-\ell)_+, \nabla(u-\ell)_+\rangle$ and
$M := \|A\nabla(u-\ell)_+\|$, we have $|M|^2 \le \Lambda E$, and
applying the abstract absorption lemma above with these choices
yields
\begin{equation*}
  \int_\Omega \eta(x)^2\, E(x)\, dx
   \le 4\,\Lambda \int_\Omega \|\nabla \eta(x)\|^2\,
     |(u-\ell)_+(x)|^2 \, dx.
\end{equation*}
Coercivity now bounds the left-hand side from below by
$\lambda \int_\Omega \eta^2 \|\nabla(u-\ell)_+\|^2$, and dividing by
$\lambda$ produces the announced inequality with constant
$4 \Lambda/\lambda$.
\end{proof}

\begin{theorem}[Ball Caccioppoli]
\label{thm:caccioppoli_weighted_on_ball_of_ballPosPart}
\lean{DeGiorgi.caccioppoli_weighted_on_ball_of_ballPosPart}
\leanok
\uses{def:MemW1pWitness, def:EllipticCoeff, def:ellipticityRatio, def:IsSubsolution, def:positivePartSub}
The same conclusion holds with $\Omega = \Ball{s}{x_0}$ and a witness
for $(u-\ell)_+$ on the ball, provided $\eta$ is smooth, nonnegative,
$|\eta| \le 1$, $\|\nabla \eta\| \le C_\eta$, and $\tsupp\,\eta
\subseteq \Ball{s}{x_0}$.
\end{theorem}

\begin{proof}
\leanok
\uses{def:MemW01p, def:deGiorgiCutoffTestGeneral, thm:deGiorgiCutoffTest_memW01p_of_truncWitness, thm:caccioppoli_weighted_of_subsolution}
See the Lean formalization.
\end{proof}

The approximation-based variant is
\leanname{caccioppoli\_weighted\_on\_ball\_of\_posPartApprox}.

\begin{theorem}[Localized energy estimate]
\label{thm:deGiorgi_energy_estimate_on_concentricBalls_of_ballPosPart}
\lean{DeGiorgi.deGiorgi_energy_estimate_on_concentricBalls_of_ballPosPart}
\leanok
\uses{def:MemW1pWitness, def:EllipticCoeff, def:ellipticityRatio, def:IsSubsolution, def:positivePartSub}
(\leanname{deGiorgi\_energy\_estimate\_on\_concentricBalls\_of\_ballPosPart})\\
Let $0 < r < s$, $A$ elliptic on $\Ball{s}{x_0}$, $u$ a subsolution.
Let $\eta$ be smooth, $0 \le \eta \le 1$ on $E$, $\eta = 1$ on
$\Ball{r}{x_0}$, $\tsupp\,\eta \subset \Ball{s}{x_0}$, and
$\|\nabla \eta\| \le C_\eta$. Then
\begin{multline*}
  \int_{\Ball{r}{x_0}} \bigl\|\nabla(u-\ell)_+(x)\bigr\|^2 \, dx \\
  \le 4\,\frac{\Lambda}{\lambda} \, C_\eta^2 \,
    \int_{\Ball{s}{x_0}} \bigl|(u-\ell)_+(x)\bigr|^2 \, dx.
\end{multline*}
\end{theorem}

\begin{proof}
\leanok
\uses{lem:EllipticCoeff_lam_nonneg, thm:caccioppoli_weighted_on_ball_of_ballPosPart}

Since $\eta = 1$ on $\Ball{r}{x_0}$, the integral on the left equals
$\int_{\Ball{r}{x_0}} \eta^2 \|\nabla(u-\ell)_+\|^2$. Monotonicity in
the integration set and the weighted Caccioppoli inequality bound
this by $4(\Lambda/\lambda)\int_{\Ball{s}{x_0}} \|\nabla \eta\|^2
|(u-\ell)_+|^2$. The pointwise bound $\|\nabla\eta\|^2 \le C_\eta^2$
on $\Ball{s}{x_0}$ then gives the claim.
\end{proof}

\section{Chebyshev bound on superlevel sets}

\begin{theorem}
\label{thm:superlevel_measureReal_le_integral_posPart}
\lean{DeGiorgi.superlevel_measureReal_le_integral_posPart}
\leanok
\uses{def:positivePartSub}
Let $(\alpha, \mu)$ be a finite measure space, $u : \alpha \to \R$,
and let $\theta < \ell$. If $|(u-\theta)_+|^2$ is integrable, then
\begin{equation*}
  \mu\bigl(\{x : \ell < u(x)\}\bigr)
   \le \frac{1}{(\ell-\theta)^2}\,
       \int |(u-\theta)_+(x)|^2 \, d\mu(x).
\end{equation*}
\end{theorem}

\begin{proof}
\leanok
\uses{lem:positivePartSub_sq_ge_gap_sq_of_lt}

Set $f(x) := |(u-\theta)_+(x)|^2$ and $\varepsilon := (\ell-\theta)^2$. The
hypothesis $\theta < \ell$ ensures $\varepsilon > 0$. By
\leanname{positivePartSub\_sq\_ge\_gap\_sq\_of\_lt}, on the superlevel set
$\{x : \ell < u(x)\}$ the value of $f$ is at least
$(\ell - \theta)^2 = \varepsilon$, so we have the set inclusion
\begin{equation*}
  \{x : \ell < u(x)\}
    \subseteq \{x : \varepsilon \le f(x)\}.
\end{equation*}
By monotonicity of measure (\leanname{measureReal\_mono}) this gives
$\mu\{\ell < u\} \le \mu\{\varepsilon \le f\}$.

Markov's inequality for nonnegative integrable functions
(\leanname{MeasureTheory.mul\_meas\_ge\_le\_integral\_of\_nonneg}) applied to
$f$ at threshold $\varepsilon$ yields
\begin{equation*}
  \varepsilon \cdot \mu\{\varepsilon \le f\}
   \le \int_{\alpha} f(x)\,d\mu(x).
\end{equation*}
Combining the set inclusion with monotone multiplication by $\varepsilon \ge 0$,
\begin{equation*}
  \varepsilon \cdot \mu\{\ell < u\}
   \le \varepsilon \cdot \mu\{\varepsilon \le f\}
   \le \int_{\alpha} f\,d\mu.
\end{equation*}
Dividing by the positive constant $\varepsilon$ produces the desired bound
$\mu\{\ell < u\} \le ((\ell - \theta)^2)^{-1}\int |(u-\theta)_+|^2\,d\mu$.
\end{proof}

\section{Pre-iteration estimate}

We now combine the energy estimate with a Sobolev/Holder
interpolation step to obtain a single-step nonlinear inequality
between truncation energies at two levels and on two concentric balls.

\begin{theorem}[Abstract pre-iteration combination]
\label{thm:deGiorgi_preiter_abstract}
\lean{DeGiorgi.deGiorgi_preiter_abstract}
\leanok
Let $d > 0$ and let $C_{\mathrm{sob}}, C_{\mathrm{en}}, I_\theta, m,
\theta < \ell$ all be nonneg numbers. Suppose
\begin{align*}
  I_\ell &\le C_{\mathrm{sob}} \cdot G \cdot m^{2/d}, \\
  m &\le (\ell - \theta)^{-2}\, I_\theta, \\
  G &\le C_{\mathrm{en}} \cdot I_\theta.
\end{align*}
Then
\begin{equation*}
  I_\ell \le C_{\mathrm{sob}}\, C_{\mathrm{en}}\,
    \bigl((\ell-\theta)^{-2} I_\theta\bigr)^{2/d}\, I_\theta.
\end{equation*}
\end{theorem}

\begin{proof}
\leanok

Raise the second inequality to the power $2/d \ge 0$, multiply by
$G$, then by $C_{\mathrm{sob}}$, using $C_{\mathrm{en}}$ to bound
$G$.
\end{proof}

\begin{theorem}[Packaged pre-iteration]
\label{thm:deGiorgi_preiter_of_energy}
\lean{DeGiorgi.deGiorgi_preiter_of_energy}
\leanok
\uses{def:positivePartSub}
With $\mu$ a finite measure and the same notation, if
\begin{itemize}
\item $I_\ell \le C_{\mathrm{sob}}\, G\, \mu(\{\ell < u\})^{2/d}$,
\item $G \le C_{\mathrm{en}} \cdot I_\theta$,
\item $\int |(u-\theta)_+|^2 \, d\mu \le I_\theta$,
\end{itemize}
then
\(
  I_\ell \le C_{\mathrm{sob}}\, C_{\mathrm{en}}\,
    ((\ell-\theta)^{-2} I_\theta)^{2/d}\, I_\theta.
\)
\end{theorem}

\begin{proof}
\leanok
\uses{thm:superlevel_measureReal_le_integral_posPart, thm:deGiorgi_preiter_abstract}

By the Chebyshev bound, $\mu(\{\ell < u\}) \le (\ell-\theta)^{-2}\,
\int|(u-\theta)_+|^2 \, d\mu \le (\ell-\theta)^{-2}\, I_\theta$. The
abstract combination lemma now applies.
\end{proof}

The cutoff Sobolev step required to discharge the Sobolev hypothesis
$I_\ell \le (C_{\mathrm{gns}})^2 \int \|\nabla \varphi\|^2\,
\mu(\{\ell < u\})^{2/d}$ is

\begin{theorem}
\label{thm:deGiorgi_cutoffSobolev_on_concentricBalls_of_ballPosPart}
\lean{DeGiorgi.deGiorgi_cutoffSobolev_on_concentricBalls_of_ballPosPart}
\leanok
\uses{def:MemW1pWitness, def:C_gns, def:positivePartSub, def:deGiorgiCutoffTestGeneral}
Under the standing hypotheses ($d > 2$, $0 < r < s$, $u$ admitting a
$W^{1,2}$ witness, $(u-\theta)_+$ admitting a $W^{1,2}$ witness on
$\Ball{s}{x_0}$, $\theta < \ell$, $\eta$ smooth with $\eta = 1$ on
$\Ball{r}{x_0}$, $|\eta| \le 1$, $\tsupp \eta \subseteq
\Ball{s}{x_0}$), there exists a $W^{1,2}$-witness $w_{\eta,\theta}$
for $\eta^2 (u-\theta)_+$ such that
\begin{multline*}
  \int_{\Ball{r}{x_0}} \bigl|(u-\ell)_+(x)\bigr|^2 \, dx
  \le (C_{\mathrm{gns}}(d,2))^2 \cdot \\
   \Bigl(\int_{\Ball{s}{x_0}} \|\nabla w_{\eta,\theta}(x)\|^2 \, dx
   \Bigr) \cdot
    \mu_s\bigl(\{\ell < u\}\bigr)^{2/d},
\end{multline*}
where $\mu_s$ denotes Lebesgue measure restricted to $\Ball{s}{x_0}$.
\end{theorem}

\begin{proof}
\leanok
\uses{def:MemW01p, lem:weak_grad_norm_memLp, lem:zero_outside_tsupport, thm:zero_extend, thm:weak_grad_unique, cor:sobolev_W01p, lem:MemW1pWitness_sub_const, lem:positivePartSub_nonneg, lem:deGiorgiCutoffTestGeneral_hasCompactSupport, lem:deGiorgiCutoffTestGeneral_tsupport_subset, thm:deGiorgiCutoffTest_memW01p_of_truncWitness}

The cutoff $\varphi := \eta^2 (u-\theta)_+$ has compact support inside
$\Ball{s}{x_0}$ and lies in $\Wop{2}(\Ball{s}{x_0})$ by
the cutoff admissibility theorem, hence its zero extension to $\R^d$
lies in $W^{1,2}(\R^d)$. The Gagliardo--Nirenberg--Sobolev embedding
gives the global $L^{2^*}$ estimate
$\|\varphi\|_{L^{2^*}} \le C_{\mathrm{gns}}(d,2) \|\nabla
\varphi\|_{L^2}$, where $2^* = 2d/(d-2)$. On $\Ball{r}{x_0}$ one has
$\eta = 1$ and $u > \ell > \theta$ implies $\varphi(x) = (u(x) -
\theta) > \ell - \theta$, so $|(u-\ell)_+|^2 \le |\varphi|^2$ on
$\{\ell < u\}$. Restricting the $L^{2^*}$ estimate to $\Ball{r}{x_0}$
and applying Holder's inequality with exponents $2^*/2$ and its dual
gives the announced bound, where the volume factor reduces to
$\mu_s(\{\ell < u\})^{2/d}$.
\end{proof}

\begin{theorem}[PDE pre-iteration on concentric balls]
\label{thm:deGiorgi_preiter_on_concentricBalls_of_ballPosPart}
\lean{DeGiorgi.deGiorgi_preiter_on_concentricBalls_of_ballPosPart}
\leanok
\uses{def:MemW1pWitness, def:C_gns, def:EllipticCoeff, def:ellipticityRatio, def:IsSubsolution, def:positivePartSub}
Under the same hypotheses ($d > 2$, $0 < r < s$, $\theta < \ell$,
$\eta$ as above), one has
\begin{multline*}
  \int_{\Ball{r}{x_0}} |(u-\ell)_+(x)|^2 \, dx
  \le (C_{\mathrm{gns}}(d,2))^2
    \cdot \bigl(8(\Lambda/\lambda + 1)\, C_\eta^2\bigr) \cdot \\
    \bigl((\ell - \theta)^{-2}
       \int_{\Ball{s}{x_0}} |(u-\theta)_+|^2 \, dx\bigr)^{2/d}
    \cdot \int_{\Ball{s}{x_0}} |(u-\theta)_+|^2 \, dx.
\end{multline*}
\end{theorem}

\begin{proof}
\leanok
\uses{thm:weak_grad_unique, lem:MemW1pWitness_mul_smooth_bounded, lem:EllipticCoeff_lam_nonneg, lem:positivePartSub_nonneg, def:deGiorgiCutoffTestGeneral, thm:caccioppoli_weighted_on_ball_of_ballPosPart, thm:deGiorgi_preiter_of_energy, thm:deGiorgi_cutoffSobolev_on_concentricBalls_of_ballPosPart}

Apply the cutoff Sobolev step with witness $w_{\eta,\theta}$, the
energy estimate $\int_{\Ball{s}{x_0}} \|\nabla w_{\eta,\theta}\|^2 \le
8(\Lambda/\lambda+1) C_\eta^2 \int_{\Ball{s}{x_0}} |(u-\theta)_+|^2$
(coming from the gradient identity for $\varphi$ together with the
weighted Caccioppoli inequality and the elementary bound
$(a+b)^2 \le 2a^2 + 2b^2$), and the packaged pre-iteration. The
identification of the gradient bound coefficient is recorded as
\leanname{deGiorgi\_cutoff\_gradient\_bound\_of\_weighted}.
\end{proof}

\section{De Giorgi iteration sequences}

We now fix the canonical scales of the iteration.

\begin{definition}
\label{def:deGiorgiTail}
\lean{DeGiorgi.deGiorgiTail, DeGiorgi.deGiorgiRadius, DeGiorgi.deGiorgiLevel, DeGiorgi.deGiorgiRecurrenceBase, DeGiorgi.deGiorgiRecurrenceCoeff}
\leanok
For $n \in \N$, $\ell^* > 0$, and $K > 0$, set
\begin{align*}
  \tau_n &:= (1/2)^{n+1}, \\
  r_n &:= (1 + \tau_n)/2, \\
  \ell_n(\ell^*) &:= \ell^* (1 - \tau_n), \\
  B_d &:= 2^{2 + 4/d}, \\
  C_d(K, \ell^*) &:= K \cdot 2^{6 + 8/d} \cdot (\ell^*)^{-4/d}.
\end{align*}
\end{definition}

\begin{lemma}
\label{lem:deGiorgiRadius_gap}
\lean{DeGiorgi.deGiorgiRadius_gap, DeGiorgi.deGiorgiLevel_gap}
\leanok
\uses{def:deGiorgiTail}
The radii $r_n$ satisfy $1/2 < r_n < 1$ and decrease to $1/2$, the
levels increase from $\ell_0 < \ell^*$ to $\ell^*$, and one has
the gap identities (\leanname{deGiorgiRadius\_gap},
\leanname{deGiorgiLevel\_gap})
\begin{equation*}
  r_n - r_{n+1} = (1/2)^{n+3}, \qquad
  \ell_{n+1}(\ell^*) - \ell_n(\ell^*) = \ell^* (1/2)^{n+2}.
\end{equation*}
\end{lemma}

\begin{proof}
\leanok
See the Lean formalization.
\end{proof}

\begin{definition}
\label{def:deGiorgiEnergySeq}
\lean{DeGiorgi.deGiorgiEnergySeq}
\leanok
\uses{def:positivePartSub, def:deGiorgiTail}
For a function $u$, base point $x_0$, and threshold $\ell^* > 0$, set
\begin{equation*}
  A_n(u, x_0, \ell^*) := \int_{\Ball{r_n}{x_0}}
     \bigl|(u - \ell_n(\ell^*))_+(x)\bigr|^2 \, dx.
\end{equation*}
\end{definition}

\section{Recurrence and abstract closeout}

\begin{theorem}[Geometric closeout for nonlinear recurrences]
\label{thm:deGiorgi_recurrence_closeout}
\lean{DeGiorgi.deGiorgi_recurrence_closeout}
\leanok
Let $Y_n \ge 0$, $C > 0$, $B > 1$, $\alpha > 0$, and assume
\begin{equation*}
  Y_{n+1} \le C\, B^n\, Y_n^{1+\alpha}, \qquad
  Y_0 \le C^{-1/\alpha}\, B^{-1/\alpha^2}.
\end{equation*}
Then $Y_n \to 0$ as $n \to \infty$.
\end{theorem}

\begin{proof}
\leanok

Set $X := C B^{1/\alpha}$, $D := X^{1/\alpha}$, $q :=
B^{-1/\alpha} < 1$, and define $W_n := B^{n/\alpha} Y_n$ and $Z_n :=
D W_n$. The hypothesis on $Y_0$ is equivalent to $Y_0 \le D^{-1}$, so
$Z_0 \le 1$. From the recurrence
\begin{equation*}
  W_{n+1} \le X \, W_n^{1+\alpha},
\end{equation*}
multiply by $D$ and use $D X = D^{1+\alpha}$ to get $Z_{n+1} \le
Z_n^{1+\alpha}$, hence by induction $Z_n \le 1$ for all $n$. This
yields $W_n \le D^{-1}$ and consequently
$Y_n \le D^{-1} q^n$, which tends to $0$ as $n \to \infty$.
\end{proof}

\begin{theorem}[Recurrence rewriting]
\label{thm:deGiorgi_preiter_to_recurrence}
\lean{DeGiorgi.deGiorgi_preiter_to_recurrence}
\leanok
\uses{def:deGiorgiTail}
If a sequence $A_n$ satisfies, for all $n$,
\begin{multline*}
  A_{n+1} \le
   \frac{K}{(r_n - r_{n+1})^2 (\ell_{n+1} - \ell_n)^{4/d}}
   \cdot A_n^{1 + 2/d},
\end{multline*}
then with $C_d(K,\ell^*)$ and $B_d$ as above one has
\begin{equation*}
  A_{n+1} \le C_d(K, \ell^*)\, B_d^n\, A_n^{1 + 2/d}.
\end{equation*}
\end{theorem}

\begin{proof}
\leanok
\uses{lem:deGiorgiRadius_gap}

By the gap identities,
$(r_n - r_{n+1})^2 = 2^{-(2n+6)}$ and
$(\ell_{n+1}-\ell_n)^{4/d} = (\ell^*)^{4/d}\, 2^{-(n+2)\cdot 4/d}$,
so the denominator equals
$(\ell^*)^{4/d}\, 2^{-(6+8/d + n(2+4/d))}$. Inverting and rearranging
yields the claimed form.
\end{proof}

\begin{theorem}[Vanishing of the canonical recurrence]
\label{thm:deGiorgi_preiter_vanishing}
\lean{DeGiorgi.deGiorgi_preiter_vanishing}
\leanok
\uses{def:deGiorgiTail}
Under $d > 2$, $K > 0$, $\ell^* > 0$, the canonical recurrence above,
and the smallness condition
\begin{equation*}
  A_0 \le C_d(K,\ell^*)^{-d/2}\, B_d^{-(d/2)^2 / 2},
\end{equation*}
the sequence $A_n$ tends to zero.
\end{theorem}

\begin{proof}
\leanok
\uses{thm:deGiorgi_recurrence_closeout, thm:deGiorgi_preiter_to_recurrence}

The pre-iteration recurrence is rewritten using
\leanname{deGiorgi\_preiter\_to\_recurrence} into the standard form with
$B := B_d$, $C := C_d(K,\ell^*)$, $\alpha := 2/d$. The smallness
hypothesis is exactly the threshold required by
\leanname{deGiorgi\_recurrence\_closeout}, so $A_n \to 0$.
\end{proof}

\section{The De Giorgi $L^\infty$ bound}

The next layer assembles the $L^\infty$ bound on the half-ball from
the recurrence machinery.

\begin{lemma}
\label{lem:ae_eq_zero_of_forall_setIntegral_sq_le_of_tendsto_zero}
\lean{DeGiorgi.ae_eq_zero_of_forall_setIntegral_sq_le_of_tendsto_zero}
\leanok
Let $f : E \to \R$ with $|f|^2 \in L^1(s)$ and let $Y_n \to 0$ be a
sequence of upper bounds for $\int_s |f|^2$. Then $f = 0$ a.e. on
$s$.
\end{lemma}

\begin{proof}
\leanok
See the Lean formalization.
\end{proof}

\begin{lemma}
\label{lem:integral_sq_halfBall_le_deGiorgiEnergySeq}
\lean{DeGiorgi.integral_sq_halfBall_le_deGiorgiEnergySeq}
\leanok
\uses{def:positivePartSub, def:deGiorgiTail, def:deGiorgiEnergySeq}
For $\ell^* > 0$, the half-ball energy is monotone in the iteration:
\begin{equation*}
  \int_{\Ball{1/2}{x_0}} |(u-\ell^*)_+|^2 \, dx
   \le A_n(u, x_0, \ell^*),
\qquad \forall n.
\end{equation*}
\end{lemma}

\begin{proof}
\leanok
\uses{lem:positivePartSub_nonneg, lem:positivePartSub_antitone}
See the Lean formalization.
\end{proof}

\begin{lemma}
\label{lem:deGiorgiEnergySeq_zero_le_integral_sq_posPartZero_on_unitBal}
\lean{DeGiorgi.deGiorgiEnergySeq_zero_le_integral_sq_posPartZero_on_unitBall}
\leanok
\uses{def:positivePartSub, def:deGiorgiTail, def:deGiorgiEnergySeq}
The first term satisfies
$A_0 \le \int_{\Bz{1}} (u_+)^2 \, dx$.
\end{lemma}

\begin{proof}
\leanok
\uses{lem:positivePartSub_nonneg, lem:positivePartSub_le_posPartZero}
See the Lean formalization.
\end{proof}

\begin{definition}
\label{def:C_DeGiorgiSmallness}
\lean{DeGiorgi.C_DeGiorgiSmallness}
\leanok
\uses{def:deGiorgiTail}
For $K > 0$ and $d \ge 1$, set
\begin{equation*}
  C_{\mathrm{small}}(d, K) :=
   \bigl(K \cdot 2^{6 + 8/d}\bigr)^{d/4}
   \cdot B_d^{d^2/8}.
\end{equation*}
\end{definition}

This is the height constant required to convert a unit-ball $L^2$
smallness assumption into the abstract initial recurrence threshold.
We collect three monotonicity/scaling identities used later:
\leanname{C\_DeGiorgiSmallness\_pos},
\leanname{C\_DeGiorgiSmallness\_mono} (monotonicity in $K$), and
\leanname{C\_DeGiorgiSmallness\_mul\_right}:
\begin{equation*}
  C_{\mathrm{small}}(d, K t) =
    C_{\mathrm{small}}(d, K) \, t^{d/4} \quad (K, t \ge 0).
\end{equation*}

\begin{definition}[Iteration constants for subsolutions]
\label{def:K_DeGiorgi_subsolution}
\lean{DeGiorgi.K_DeGiorgi_subsolution, DeGiorgi.C_DeGiorgi_subsolution}
\leanok
\uses{def:C_gns, lem:smoothTransition_nnnorm_deriv_bounded, def:EllipticCoeff, def:ellipticityRatio, def:C_DeGiorgiSmallness}
For an elliptic coefficient field $A$ on $\Bz{1}$ set
\begin{equation*}
  K_{\mathrm{DG}}(d, A) := (C_{\mathrm{gns}}(d, 2))^2 \cdot
    8\bigl(\Lambda/\lambda + 1\bigr) \cdot (2 M_{\mathrm{st}})^2 + 1,
\end{equation*}
\begin{equation*}
  C_{\mathrm{DG}}(d, A) := C_{\mathrm{small}}(d, K_{\mathrm{DG}}(d, A)).
\end{equation*}
Both are positive (\leanname{K\_DeGiorgi\_subsolution\_pos},
\leanname{C\_DeGiorgi\_subsolution\_pos}). The constant
$M_{\mathrm{st}}$ is the standard cutoff gradient bound from the
Stampacchia layer.
\end{definition}

\begin{theorem}[Iteration package]
\label{thm:deGiorgi_iteration_package_on_unitBall_of_subsolution}
\lean{DeGiorgi.deGiorgi_iteration_package_on_unitBall_of_subsolution}
\leanok
\uses{def:EllipticCoeff, def:IsSubsolution, def:positivePartSub, def:deGiorgiTail, def:deGiorgiEnergySeq, def:K_DeGiorgi_subsolution}
Let $d > 2$, $A$ elliptic on $\Bz{1}$, and $u$ a subsolution. Then
$K_{\mathrm{DG}}(d,A) > 0$ and for every $\ell^* > 0$:
\begin{enumerate}
\item $|(u-\ell^*)_+|^2$ is integrable on $\Ball{1/2}{0}$;
\item for every $n \in \N$, $|(u-\ell_n(\ell^*))_+|^2$ is integrable
on $\Ball{r_n}{0}$;
\item the canonical pre-iteration recurrence holds:
\begin{multline*}
  A_{n+1}(u,0,\ell^*) \le \\
   \frac{K_{\mathrm{DG}}(d,A)}{(r_n - r_{n+1})^2 \,
     (\ell_{n+1}(\ell^*)-\ell_n(\ell^*))^{4/d}} \cdot
   A_n(u,0,\ell^*)^{1 + 2/d}.
\end{multline*}
\end{enumerate}
\end{theorem}

\begin{proof}
\leanok
\uses{def:MemW1p, def:MemW1pWitness, def:MemW01p, lem:weak_grad_norm_memLp, lem:zero_outside_tsupport, thm:zero_extend, thm:weak_grad_unique, def:C_gns, cor:sobolev_W01p, def:CampanatoBall, def:HasCampanatoBound, lem:smoothTransition_nnnorm_deriv_bounded, thm:Cutoff_exists, thm:w12-approx, lem:MemW1pWitness_sub_const, lem:MemW1pWitness_mul_smooth_bounded, thm:MemW1pWitness_posPart, def:ellipticityRatio, lem:EllipticCoeff_lam_nonneg, lem:positivePartSub_nonneg, def:deGiorgiCutoffTestGeneral, lem:deGiorgiCutoffTestGeneral_hasCompactSupport, lem:deGiorgiCutoffTestGeneral_tsupport_subset, thm:deGiorgiCutoffTest_memW01p_of_truncWitness, def:positivePartSub_memW1pWitness_on_ball, thm:caccioppoli_weighted_of_subsolution, thm:deGiorgi_preiter_of_energy, lem:deGiorgiRadius_gap}

Take a $W^{1,2}(\Bz{1})$ witness for $u$ from
$\leanname{IsSubsolution.memW1p}$. The smooth approximation package
on the unit ball produces, for any shift $\theta$, an approximating
sequence $\psi_n \in C^1_c$ with $\psi_n \to u - \theta$ in $L^2$ and
$\nabla \psi_n \to \nabla u$ in $L^2$. From this one constructs
witnesses for $(u-\ell^*)_+$ and $(u-\ell_n(\ell^*))_+$ on $\Bz{1}$
via \leanname{positivePartSub\_memW1pWitness\_on\_ball}, and the
required integrability follows by restriction to $\Ball{r_n}{0}$. For
the recurrence, fix $n$ and apply the PDE pre-iteration theorem
\leanname{deGiorgi\_preiter\_on\_concentricBalls\_of\_ballPosPart} to $r =
r_{n+1}$, $s = r_n$, $\theta = \ell_n(\ell^*)$, $\ell =
\ell_{n+1}(\ell^*)$, with a smooth Stampacchia cutoff $\eta_n$
satisfying $\eta_n = 1$ on $\Ball{r_{n+1}}{0}$, $\tsupp\,\eta_n
\subset \Ball{r_n}{0}$, and $\|\nabla \eta_n\| \le 2 M_{\mathrm{st}}
/ (r_n - r_{n+1})$, so that the constant becomes
$(C_{\mathrm{gns}}(d,2))^2 \cdot 8(\Lambda/\lambda + 1)
(2 M_{\mathrm{st}})^2 \le K_{\mathrm{DG}}(d,A)$.
\end{proof}

\begin{theorem}[Vanishing version]
\label{thm:linfty_subsolution_DeGiorgi_ae_zero_of_smallness}
\lean{DeGiorgi.linfty_subsolution_DeGiorgi_ae_zero_of_smallness}
\leanok
\uses{def:positivePartSub, def:deGiorgiTail, def:deGiorgiEnergySeq}
With $d > 2$, $K, \ell^* > 0$, the integrability assumptions on the
truncations, the canonical recurrence, and the abstract smallness
condition
\begin{equation*}
  A_0(u,x_0,\ell^*) \le
    C_d(K,\ell^*)^{-d/2}\, B_d^{-(d/2)^2 / 2},
\end{equation*}
one has $(u-\ell^*)_+ = 0$ almost everywhere on $\Ball{1/2}{x_0}$.
\end{theorem}

\begin{proof}
\leanok
\uses{thm:deGiorgi_preiter_vanishing, lem:ae_eq_zero_of_forall_setIntegral_sq_le_of_tendsto_zero, lem:integral_sq_halfBall_le_deGiorgiEnergySeq}

By \leanname{deGiorgi\_preiter\_vanishing}, $A_n \to 0$. Combined with
the half-ball monotonicity lemma, this gives
$\int_{\Ball{1/2}{x_0}} |(u-\ell^*)_+|^2 = 0$, whence
$(u-\ell^*)_+ = 0$ a.e.\ on the half-ball by the abstract zero
lemma.
\end{proof}

\begin{corollary}
\label{cor:linfty_subsolution_DeGiorgi_ae_of_smallness}
\lean{DeGiorgi.linfty_subsolution_DeGiorgi_ae_of_smallness}
\leanok
\uses{def:positivePartSub, def:deGiorgiTail, def:deGiorgiEnergySeq}
Under the same hypotheses, $\max(u, 0) \le \ell^*$ a.e.\ on
$\Ball{1/2}{x_0}$.
\end{corollary}

\begin{proof}
\leanok
\uses{lem:le_of_positivePartSub_eq_zero, thm:linfty_subsolution_DeGiorgi_ae_zero_of_smallness}

$(u-\ell^*)_+(x) = 0$ implies $u(x) \le \ell^*$ by
\leanname{le\_of\_positivePartSub\_eq\_zero}, and $0 \le \ell^*$.
\end{proof}

\begin{theorem}[Unit-ball smallness packaging]
\label{thm:linfty_subsolution_DeGiorgi_ae_of_unitBall_initial_energy_sm}
\lean{DeGiorgi.linfty_subsolution_DeGiorgi_ae_of_unitBall_initial_energy_smallness}
\leanok
\uses{def:positivePartSub, def:deGiorgiTail, def:deGiorgiEnergySeq}
Replacing the abstract smallness on $A_0$ by the same threshold for
$\int_{\Bz{1}} (u_+)^2$, the same conclusion holds.
\end{theorem}

\begin{proof}
\leanok
\uses{lem:deGiorgiEnergySeq_zero_le_integral_sq_posPartZero_on_unitBal, cor:linfty_subsolution_DeGiorgi_ae_of_smallness}

The bound $A_0 \le \int_{\Bz{1}} (u_+)^2$ converts the assumption
into the previous form.
\end{proof}

\begin{theorem}[Height-threshold form]
\label{thm:linfty_subsolution_DeGiorgi_ae_of_unitBall_height_bound}
\lean{DeGiorgi.linfty_subsolution_DeGiorgi_ae_of_unitBall_height_bound}
\leanok
\uses{def:positivePartSub, def:deGiorgiTail, def:deGiorgiEnergySeq, def:C_DeGiorgiSmallness}
Under the same setup, if
\begin{equation*}
  C_{\mathrm{small}}(d,K)\,
    \Bigl(\int_{\Bz{1}} (u_+)^2\Bigr)^{1/2} \le \ell^*,
\end{equation*}
then $\max(u,0) \le \ell^*$ a.e.\ on $\Ball{1/2}{x_0}$.
\end{theorem}

\begin{proof}
\leanok
\uses{thm:linfty_subsolution_DeGiorgi_ae_of_unitBall_initial_energy_sm}

Squaring and using the formula for $C_{\mathrm{small}}$ shows that
the height-bound implies the abstract smallness condition above.
This step is the lemma
\leanname{deGiorgi\_initial\_energy\_small\_of\_height\_bound}.
\end{proof}

\begin{theorem}[De Giorgi $L^\infty$ bound]
\label{thm:linfty_subsolution_DeGiorgi}
\lean{DeGiorgi.linfty_subsolution_DeGiorgi}
\leanok
\uses{def:EllipticCoeff, def:IsSubsolution, def:K_DeGiorgi_subsolution}
Let $d > 2$, $A$ elliptic on $\Bz{1}$, and $u$ a subsolution with
$|\max(u,0)|^2 \in L^1(\Bz{1})$. Then
\begin{equation*}
  \max(u(x), 0) \le C_{\mathrm{DG}}(d, A) \cdot
    \Bigl(\int_{\Bz{1}} |\max(u,0)|^2\, dx\Bigr)^{1/2}
\end{equation*}
for almost every $x \in \Ball{1/2}{0}$.
\end{theorem}

\begin{proof}
\leanok
\uses{def:deGiorgiTail, lem:ae_eq_zero_of_forall_setIntegral_sq_le_of_tendsto_zero, def:C_DeGiorgiSmallness, thm:deGiorgi_iteration_package_on_unitBall_of_subsolution, thm:linfty_subsolution_DeGiorgi_ae_of_unitBall_height_bound}

Set $I_0 := \int_{\Bz{1}} |u_+|^2 \, dx$. If $I_0 = 0$ then by
monotonicity $|u_+|^2 = 0$ a.e.\ on $\Ball{1/2}{0}$, so the
inequality holds trivially. Otherwise let $\ell^* := C_{\mathrm{DG}}
(d,A) \sqrt{I_0} > 0$. The iteration package supplies the integrability
of all canonical truncations and the canonical recurrence. The
height-threshold theorem then yields $\max(u,0) \le \ell^* =
C_{\mathrm{DG}}(d,A) \sqrt{I_0}$ a.e.\ on $\Ball{1/2}{0}$.
\end{proof}

\begin{corollary}[Existence of a level threshold]
\label{cor:linfty_subsolution_DeGiorgi_exists_threshold}
\lean{DeGiorgi.linfty_subsolution_DeGiorgi_exists_threshold}
\leanok
\uses{def:EllipticCoeff, def:IsSubsolution, def:K_DeGiorgi_subsolution}
Under the same hypotheses there exists $\ell^* > 0$ with
$C_{\mathrm{DG}}(d,A) \sqrt{I_0} \le \ell^*$ such that
$\max(u,0) \le \ell^*$ a.e.\ on $\Ball{1/2}{0}$.
\end{corollary}

\begin{proof}
\leanok
\uses{def:C_gns, lem:smoothTransition_nnnorm_deriv_bounded, def:ellipticityRatio, lem:EllipticCoeff_lam_nonneg, def:deGiorgiTail, def:C_DeGiorgiSmallness, thm:linfty_subsolution_DeGiorgi}
See the Lean formalization.
\end{proof}

\section{Normalised constants}

\begin{definition}
\label{def:K_DeGiorgi_subsolution_normalized}
\lean{DeGiorgi.K_DeGiorgi_subsolution_normalized, DeGiorgi.C_DeGiorgi_subsolution_normalized}
\leanok
\uses{def:C_gns, lem:smoothTransition_nnnorm_deriv_bounded, def:C_DeGiorgiSmallness}
The dimension-only counterparts are
\begin{equation*}
  K^{\mathrm{norm}}_{\mathrm{DG}}(d) :=
    2 \cdot (C_{\mathrm{gns}}(d,2))^2 \cdot 8\, (2 M_{\mathrm{st}})^2 + 1,
\end{equation*}
\begin{equation*}
  C^{\mathrm{norm}}_{\mathrm{DG}}(d) :=
    C_{\mathrm{small}}(d, K^{\mathrm{norm}}_{\mathrm{DG}}(d)).
\end{equation*}
\end{definition}

\begin{lemma}
\label{lem:C_DeGiorgi_subsolution_le_normalized}
\lean{DeGiorgi.C_DeGiorgi_subsolution_le_normalized}
\leanok
\uses{def:ellipticityRatio, def:K_DeGiorgi_subsolution, def:K_DeGiorgi_subsolution_normalized}
For a normalised elliptic coefficient field $A$ on $\Bz{1}$ (i.e.\
$\lambda = 1$, so $\Lambda \ge 1$ and the ratio
$\Lambda/\lambda = \Lambda$),
\begin{equation*}
  C_{\mathrm{DG}}(d, A) \le
    C^{\mathrm{norm}}_{\mathrm{DG}}(d) \cdot \Lambda^{d/4}.
\end{equation*}
\end{lemma}

\begin{proof}
\leanok
\uses{def:C_gns, lem:smoothTransition_nnnorm_deriv_bounded, def:EllipticCoeff, lem:EllipticCoeff_lam_nonneg, def:deGiorgiTail, def:C_DeGiorgiSmallness}

Comparing the formulas, $K_{\mathrm{DG}}(d,A) = c (\Lambda+1) + 1$ and
$K^{\mathrm{norm}}_{\mathrm{DG}}(d) = 2c + 1$ where $c =
(C_{\mathrm{gns}}(d,2))^2 \cdot 8 (2 M_{\mathrm{st}})^2$. Since
$\Lambda \ge 1$ one has $c(\Lambda + 1) + 1 \le (2c+1)\Lambda$,
i.e.\ $K_{\mathrm{DG}}(d,A) \le K^{\mathrm{norm}}_{\mathrm{DG}}(d) \cdot
\Lambda$. Monotonicity of $C_{\mathrm{small}}$ in $K$ followed by the
multiplicative scaling identity
$C_{\mathrm{small}}(d, K \cdot \Lambda) = C_{\mathrm{small}}(d, K)\,
\Lambda^{d/4}$ gives the claim.
\end{proof}

\begin{theorem}[Normalised De Giorgi $L^\infty$ bound]
\label{thm:linfty_subsolution_DeGiorgi_normalized}
\lean{DeGiorgi.linfty_subsolution_DeGiorgi_normalized}
\leanok
\uses{def:ellipticityRatio, def:IsSubsolution, def:K_DeGiorgi_subsolution_normalized}
Let $d > 2$, $A$ a normalised elliptic coefficient field on
$\Bz{1}$, and $u$ a subsolution with
$|\max(u,0)|^2 \in L^1(\Bz{1})$. Then
\begin{equation*}
  \max(u(x), 0) \le
    C^{\mathrm{norm}}_{\mathrm{DG}}(d) \cdot \Lambda^{d/4} \cdot
    \Bigl(\int_{\Bz{1}} |\max(u,0)|^2 \, dx\Bigr)^{1/2}
\end{equation*}
for almost every $x \in \Ball{1/2}{0}$.
\end{theorem}

\begin{proof}
\leanok
\uses{def:K_DeGiorgi_subsolution, thm:linfty_subsolution_DeGiorgi, lem:C_DeGiorgi_subsolution_le_normalized}

Combine \leanname{linfty\_subsolution\_DeGiorgi} with the previous
lemma comparing $C_{\mathrm{DG}}(d,A)$ to
$C^{\mathrm{norm}}_{\mathrm{DG}}(d) \, \Lambda^{d/4}$, multiplying
by the nonnegative factor $\sqrt{I_0}$.
\end{proof}

This concludes the formal De Giorgi $L^2 \to L^\infty$ chapter. The
two top-level interfaces \leanname{DeGiorgiIteration} and
\leanname{DeGiorgiTheory} simply re-export the constructions and
theorems above so that downstream chapters (weak Harnack, Harnack,
Holder) can rely on the De Giorgi bound directly.


%% ========================================================================
